\documentclass{umk-matematika}

\begin{document}

\maketitlepage
    {Практическая работа №17}
    {Показательные уравнения}
    {}

\section{Фундамент (1--4 балла)}

\textbf{Задача 1 (1 балл).} Решите уравнение: $2^x = 8$

\textbf{Задача 2 (2 балла).} Решите уравнение: $5^x = 25$

\textbf{Задача 3 (3 балла).} Решите уравнение: $3^x = 1$

\textbf{Задача 4 (4 балла).} Решите уравнение: $10^x = 1000$

\section{Стены (5--7 баллов)}

\textbf{Задача 5 (5 баллов).} Решите уравнение: $5^{2x-1} = 125$

\textbf{Задача 6 (6 баллов).} Решите уравнение: $3^{x+2} + 3^x = 90$

\textbf{Задача 7 (7 баллов).} Решите уравнение: $2^{2x} - 6 \cdot 2^x + 8 = 0$

\section{Кровля (8--10 баллов)}

\textbf{Задача 8 (8 баллов).} Решите уравнение: $4^x - 3 \cdot 2^{x+1} + 8 = 0$

\textbf{Задача 9 (9 баллов).} Решите уравнение: $9^x - 4 \cdot 3^x + 3 = 0$

\textbf{Задача 10 (10 баллов).} Население города растёт по закону $N(t) = N_0 \cdot 2^{\frac{t}{10}}$, где t — годы. Через сколько лет население удвоится?

\newpage
\section*{Ответы}

\begin{enumerate}
  \item $x = 3$
  \item $x = 2$
  \item $x = 0$
  \item $x = 3$
  \item $x = 2$
  \item $x = 2$
  \item $x = 1;\ x = 2$
  \item $x = 1;\ x = 2$
  \item $x = 0;\ x = 1$
  \item через 10 лет.
\end{enumerate}

\end{document}